\documentclass{book}

\def\title{Elementary Analysis}
\def\author{Nathanael Chwojko-Srawley}


%  ╔══════════════════════════════════════════════════════════╗
%  ║                         packages                         ║
%  ╚══════════════════════════════════════════════════════════╝

\usepackage[a4paper, total={6in, 8in}]{geometry}
\usepackage[answerdelayed]{exercise}
\usepackage[font=itshape]{quoting}
\usepackage[most]{tcolorbox}       % Makes nice boxes for thms, cor, prop, or anything else.
\usepackage{amsmath, amssymb, amsthm} % Used to write better mathematical expressions (ex. \mathbb{R})
\usepackage{array}            % \begin{table}{>{\em}c c} -> 1st column italic (bfseries for bold)
\usepackage{bbm}                     % for mathbbm{1}
\usepackage{blindtext}
\usepackage{commutative-diagrams}
\usepackage{dsfont}
\usepackage{dynkin-diagrams}   % for Dynkin diagrams
\usepackage{enumitem}
\usepackage{etoolbox}          % for Dynkin diagrams
\usepackage{euscript}
\usepackage{extarrows}               % Used to compile any UTF-8 Character (and supports Unicode)
\usepackage{fancyhdr}
\usepackage{float}            % package with ``H'' for figures and tables, ex. Images, tables, etc
\usepackage{graphicx}         % Let's you use images in LaTeX; ex the \includegraphics command.
\usepackage{hhline}
\usepackage{hyperref}
\usepackage{cleveref}
\usepackage{inputenc}         % Used to compile any UTF-8 Character (and supports Unicode)
\usepackage{listings}         % Create nice colour code blocks (customized in preamble)
\usepackage{longtable}        % create tables that extend past one page
\usepackage{makeidx}
\usepackage{mathrsfs}
\usepackage{mathtools}
\usepackage{microtype}        % makes nice text. Fixes some under-/over- filled box problems.
\usepackage{multicol}         % For multiple columns
\usepackage{pgfplots}
\usepackage{scalerel}
\usepackage{setspace}         % More flexibility on spacing in document.
\usepackage{stackengine}
\usepackage{stackrel}         % for left-right arrows, staking symbol above and bellow
\usepackage{thmtools}         % Customize theorem environment (in preamble).
\usepackage{tikz-cd}          % for commutative diagrams https://tex.stackexchange.com/questions/419221/how-to-do-the-pushout-with-universal-property
\usepackage{tikz}                    % for commutative diagrams https://tex.stackexchange.com/questions/419221/how-to-do-the-pushout-with-universal-property
\usepackage{titlesec}
\usepackage{wasysym}          % emoji's!
\usepackage{wrapfig}          % To wrap text around figure
\usepackage{xcolor}           % Colour text.
\usepackage{subcaption}

% \usepackage{comment}

% \usepackage{CJKutf8} % for some chinese

% \usepackage{dutchcal} % for lower case mathcal
% \pdfsuppresswarningpagegroup=1

% ╔═════════════════════════════════════════════════════╗
% ║           for figures via inkspace                  ║
% ╚═════════════════════════════════════════════════════╝

% to import figures via: https://github.com/gillescastel/inkscape-figures
% \usepackage{import}
% \usepackage{pdfpages}
% \usepackage{transparent}
%
%
% % This command is the ONLY command imported via the packages snippet, think of it as a tiny package written out
% \newcommand{\incfig}[2][1]{%
%     \def\svgwidth{#1\columnwidth}
%     \import{./figures/}{#2.pdf_tex}
% }

% ╔═══════════════════════════════════════════════════════════╗
% ║                 bibliography and citing                   ║
% ╚═══════════════════════════════════════════════════════════╝

\usepackage{csquotes}
% \usepackage[backend=biber,citestyle=alphabetic,bibstyle=authortitle]{biblatex}
\usepackage[backend=biber]{biblatex}

\addbibresource{bibliography.bib}

% ╔═══════════════════════════════════════════════════════════╗
% ║                    colours (colors)                       ║
% ╚═══════════════════════════════════════════════════════════╝

\definecolor{dkgreen}{rgb}{0,0.6,0}
\definecolor{gray}{rgb}{0.5,0.5,0.5}
\definecolor{mauve}{rgb}{0.58,0,0.82}
\definecolor{lightyellow}{rgb}{1,1,0.6}
\definecolor{reddish}{HTML}{A01A3A}

%  ╔══════════════════════════════════════════════════════════╗
%  ║                       book format                        ║
%  ╚══════════════════════════════════════════════════════════╝

%% check if using xelatex rather than pdflatex
\usepackage[T1]{fontenc}
\usepackage{lmodern}
\ifpdf
\usepackage{pdfcolmk}
\fi
\ifxetex
\usepackage{fontspec}
\fi
\usepackage{pifont}

\providecommand{\HUGE}{\Huge}% if not using memoir
\newlength{\drop}% for my convenience
%% select a (FontSite) font by its font family ID
\newcommand*{\FSfont}[1]{\fontencoding{T1}\fontfamily{#1}\selectfont}
%% if you don’t have the FontSite fonts either \renewcommand*{\FSfont}[1]{}
%% or use your own choice of family.
%% select a (TeX Font) font by its font family ID
\newcommand*{\TXfont}[1]{\fontencoding{T1}\fontfamily{#1}\selectfont}
%% Generic publisher’s logo
\newcommand*{\plogo}{\fbox{$\mathcal{PL}$}}

%% Some shades
\definecolor{Dark}{gray}{0.2}
\definecolor{MedDark}{gray}{0.4}
\definecolor{Medium}{gray}{0.6}
\definecolor{Light}{gray}{0.8}
%%%% Additional font series macros
\makeatletter
%%%% light series
%% e.g., kernel doc, section s: line 12 or thereabouts
\DeclareRobustCommand\ltseries
{\not@math@alphabet\ltseries\relax
\fontseries\ltdefault\selectfont}
%% e.g., kernel doc, section t: line 32 or thereabouts
\newcommand{\ltdefault}{l}
%% e.g., kernel doc, section v: line 19 or thereabouts
\DeclareTextFontCommand{\textlt}{\ltseries}
% heavy(bold) series
\DeclareRobustCommand\hbseries
{\not@math@alphabet\hbseries\relax
\fontseries\hbdefault\selectfont}
\newcommand{\hbdefault}{hb}
\DeclareTextFontCommand{\texthb}{\hbseries}
\makeatother


\newcommand{\titleUL}{\begingroup% University of Liege
\drop=0.1\textheight
\vspace*{\drop}
\begin{center}
% {\LARGE\textsc{University of Toronto}}\\[\drop]
% University logo
% {\LARGE \plogo}\\[\drop]
{\LARGE {\quad}}\\[\drop]

\rule{\textwidth}{1pt}\par
\vspace{0.5\baselineskip}
{\huge\bfseries \title\\[0.2\baselineskip]
\large --- \emph{EYNTKA} Series ---}\\[0.5\baselineskip]
\rule{\textwidth}{1pt}\par
\vfill
{\Large\textsc{\author}}
\vfill
\vfill
{\large \today}
\end{center}
\endgroup}

%header information
\pagestyle{fancy}
\setlength\headheight{20pt}
\renewcommand{\sectionmark}[1]{\markright{\thesection.\ #1}}
\renewcommand{\chaptermark}[1]{\markboth{#1}{}}
\fancyfoot[C,CO]{\textcolor{gray}{Nathanael Chwojko-Srawley}}
\fancyfoot[R,RO]{\thepage}{}

\titleformat
{\chapter} % command
[display] % shape
{\bfseries\Huge\itshape} % format
{\thechapter} % label
{0.5ex} % sep
{
    \rule{\textwidth}{1pt}
    \vspace{2ex}
    \centering
} % before-code
[
\vspace{-0.5ex}%
\rule{\textwidth}{1pt}
] % after-code

% Create a new page style for the first page
\fancypagestyle{firstpage}{
  \fancyhf{} % Clear header and footer
  \renewcommand{\headrulewidth}{0pt} % This removes the header line
  \fancyfoot[C]{\thepage}
}

% Apply the first page style conditionally
\AtBeginDocument{\thispagestyle{firstpage}}

\renewcommand{\maketitle}{\titleUL}

%  ╔══════════════════════════════════════════════════════════╗
%  ║                     Spacing Format                       ║
%  ╚══════════════════════════════════════════════════════════╝

\setlength{\parindent}{0pt} % don't like indentation infront of paragraphs
\setlength{\parskip}{1ex}   %so that there is still space between paragarphs


%  ╔══════════════════════════════════════════════════════════╗
%  ║                       Shortcuts                          ║
%  ╚══════════════════════════════════════════════════════════╝

%\ExerciseLevelInToc

% I like original mathcal better, but need lowercase.
\DeclareFontFamily{OT1}{pzc}{}
\DeclareFontShape{OT1}{pzc}{m}{it}{<-> s * [1.10] pzcmi7t}{}
\DeclareMathAlphabet{\mathpzc}{OT1}{pzc}{m}{it}


% mathbb
\newcommand{\A}{{\mathbb{A}}}
\newcommand{\C}{{\mathbb{C}}}
\newcommand{\F}{{\mathbb{F}}}
\newcommand{\T}{{\mathbb{T}}}
\newcommand{\N}{{\mathbb{N}}}
\newcommand{\Q}{{\mathbb{Q}}}
\newcommand{\R}{{\mathbb{R}}}
\newcommand{\V}{{\mathbb{V}}}
\newcommand{\Z}{{\mathbb{Z}}}
\newcommand{\BI}{{\mathbb{I}}}
\renewcommand{\H}{{\mathbb{H}}}
\renewcommand{\P}{{\mathbb{P}}}


% mathcal
% \newcommand{\co}{{\mathfrak{o}}} %mathcal{o} looks like wreath product, so I changed it to mathfrak
\newcommand{\cc}{{\mathpzc{c}}}
\newcommand{\co}{{\mathpzc{o}}}
\newcommand{\cp}{{\mathpzc{p}}}
\newcommand{\CA}{{\mathcal{A}}}
\newcommand{\CB}{{\mathcal{B}}}
\newcommand{\CC}{{\mathcal{C}}}
\newcommand{\CD}{{\mathcal{D}}}
\newcommand{\CF}{{\mathcal{F}}}
\newcommand{\CG}{{\mathcal{G}}}
\newcommand{\CH}{{\mathcal{H}}}
\newcommand{\CI}{{\mathcal{I}}}
\newcommand{\CK}{{\mathcal{K}}}
\newcommand{\CL}{{\mathcal{L}}}
\newcommand{\CN}{{\mathcal{N}}}
\newcommand{\CM}{{\mathcal{M}}}
\newcommand{\CO}{{\mathcal{O}}}
\newcommand{\CU}{{\mathcal{U}}}
\newcommand{\CQ}{{\mathcal{Q}}}
\newcommand{\CR}{{\mathcal{R}}}
\newcommand{\CS}{{\mathcal{S}}}
\newcommand{\CT}{{\mathcal{T}}}
\newcommand{\CV}{{\mathcal{V}}}
\newcommand{\CZ}{{\mathcal{Z}}}
% EXCPETION:
\newcommand{\CP}{{\mathcal{P}}}
% \newcommand{\CP}{\mathbb{C}\mathrm{P}}
\newcommand{\RP}{\mathbb{R}\mathrm{P}}


%Euscript
\newcommand{\EA}{{\EuScript{A}}}
\newcommand{\EB}{{\EuScript{B}}}
\newcommand{\EC}{{\EuScript{C}}}
\newcommand{\ED}{{\EuScript{D}}}
\newcommand{\EF}{{\EuScript{F}}}
\newcommand{\EL}{{\EuScript{L}}}
\newcommand{\EO}{{\EuScript{O}}}
\newcommand{\ET}{{\EuScript{T}}}


%mathfrak (lowercase)
\newcommand{\fa}{{\mathfrak{a}}}
\newcommand{\fb}{{\mathfrak{b}}}
\newcommand{\fc}{{\mathfrak{c}}}
\newcommand{\fd}{{\mathfrak{d}}}
\newcommand{\fm}{{\mathfrak{m}}}
\newcommand{\fn}{{\mathfrak{n}}}
\newcommand{\fo}{{\mathfrak{o}}}
\newcommand{\fp}{{\mathfrak{p}}}
\newcommand{\fq}{{\mathfrak{q}}}


%mathfrak (upper case)
\newcommand{\FA}{{\mathfrak{A}}}
\newcommand{\FB}{{\mathfrak{B}}}
\newcommand{\FC}{{\mathfrak{C}}}
\newcommand{\FD}{{\mathfrak{D}}}
\newcommand{\FK}{{\mathfrak{K}}}
\newcommand{\FM}{{\mathfrak{M}}}
\newcommand{\FN}{{\mathfrak{N}}}
\newcommand{\FO}{{\mathfrak{O}}}
\newcommand{\FP}{{\mathfrak{P}}}
\newcommand{\FX}{{\mathfrak{X}}}


%mathscr
\newcommand{\ff}{{\mathscr{F}}}
\newcommand{\SA}{\mathscr{A}}
\newcommand{\SC}{\mathscr{C}}
\newcommand{\SF}{\mathscr{F}}
\newcommand{\SG}{\mathscr{G}}
\newcommand{\SH}{\mathscr{H}}
\newcommand{\SI}{\mathscr{I}}
\newcommand{\SK}{\mathscr{K}}
\newcommand{\SN}{\mathscr{N}}
\newcommand{\SQ}{\mathscr{Q}}
\newcommand{\SR}{\mathscr{R}}
\newcommand{\SV}{\mathscr{V}}
\newcommand{\SZ}{\mathscr{Z}}


% operatorname -- 2 letters (lowercase and upper case)
\newcommand{\ab}{{\operatorname{ab}}}
\newcommand{\ad}{{\operatorname{ad}}}
\newcommand{\ch}{{\operatorname{ch}}}
\newcommand{\ev}{{\operatorname{ev}}}
\newcommand{\id}{{\operatorname{id}}}
\newcommand{\im}{{\operatorname{im}}}
\newcommand{\nm}{{\operatorname{Nm}}}
\newcommand{\op}{{\operatorname{op}}}
\newcommand{\rk}{{\operatorname{rk}}}
\newcommand{\sh}{{\operatorname{sh}}}
\newcommand{\tr}{{\operatorname{tr}}}

\newcommand{\Ab}{{\operatorname{Ab}}}
\newcommand{\Br}{{\operatorname{Br}}}
\newcommand{\Ch}{{\operatorname{Ch}}}
\newcommand{\Cl}{{\operatorname{Cl}}}
\newcommand{\Nm}{{\operatorname{Nm}}}
\newcommand{\SL}{{\operatorname{SL}}}
\newcommand{\SO}{{\operatorname{SO}}}
\newcommand{\Sp}{{\operatorname{Sp}}}
\newcommand{\GL}{{\operatorname{GL}}}
\newcommand{\PGL}{{\operatorname{PGL}}}

\renewcommand{\Re}{{\operatorname{Re}}}
\renewcommand{\Im}{{\operatorname{Im}}}


%categories
\newcommand{\Grp}{{\textbf{Grp}}}
\newcommand{\Fld}{{\textbf{Fld}}}
\newcommand{\Sh}{{\textbf{Sh}}}
\newcommand{\Set}{{\textbf{Set}}}
\newcommand{\Mod}{{\textbf{Mod}}}
\newcommand{\PSh}{{\textbf{PSh}}}
\newcommand{\Sch}{{\textbf{Sch}}}
\newcommand{\Top}{{\textbf{Top}}}
\newcommand{\RgSp}{{\textbf{RgSp}}}
\newcommand{\Ring}{{\textbf{Ring}}}


%operatorname -- 3+ letters (lowercase and upper case)
\newcommand{\rad}{{\operatorname{rad}}}
\newcommand{\sep}{{\operatorname{sep}}}
\newcommand{\res}{{\operatorname{res}}}
\newcommand{\sgn}{{\operatorname{sgn}}}
\newcommand{\pre}{{\operatorname{pre}}}
\newcommand{\ord}{{\operatorname{ord}}}
\newcommand{\vol}{{\operatorname{vol}}}
\newcommand{\lcm}{{\operatorname{lcm}}}

\newcommand{\conj}{{\operatorname{conj}}}
\newcommand{\supp}{{\operatorname{supp}}}
\newcommand{\rank}{{\operatorname{rank}}}
\newcommand{\disc}{{\operatorname{disc}}}
\newcommand{\stab}{{\operatorname{stab}}}
\newcommand{\kdim}{{\operatorname{kdim}}}
\newcommand{\diag}{{\operatorname{diag}}}

\newcommand{\trdeg}{{\operatorname{trdeg}}}
\newcommand{\coker}{{\operatorname{coker}}}
\newcommand{\codim}{{\operatorname{codim}}}

\newcommand{\Ann}{{\operatorname{Ann}}}
\newcommand{\Aut}{{\operatorname{Aut}}}
\newcommand{\Der}{{\operatorname{Der}}}
\newcommand{\Emb}{{\operatorname{Emb}}}
\newcommand{\End}{{\operatorname{End}}}
\newcommand{\Ext}{{\operatorname{Ext}}}
\newcommand{\Gal}{{\operatorname{Gal}}}
\newcommand{\Hom}{{\operatorname{Hom}}}
\newcommand{\Ind}{{\operatorname{Ind}}}
\newcommand{\Inn}{{\operatorname{Inn}}}
\newcommand{\Int}{{\operatorname{int}}}
\newcommand{\Irr}{{\operatorname{Irr}}}
\newcommand{\Jac}{{\operatorname{Jac}}}
\newcommand{\Mor}[1]{\operatorname{Mor}(\textbf{#1})}
\newcommand{\Nil}{{\operatorname{Nil}}}
\newcommand{\Obj}{{\operatorname{Obj}}}
\newcommand{\Out}{{\operatorname{Out}}}
\newcommand{\Pic}{{\operatorname{Pic}\,}}
\newcommand{\Rep}{{\operatorname{Rep}}}
\newcommand{\Res}{{\operatorname{Res}}}
\newcommand{\Spl}{{\operatorname{Spl}}}
\newcommand{\Syl}{{\operatorname{Syl}}}
\newcommand{\Sym}{{\operatorname{Sym}}}
\newcommand{\Tor}{{\operatorname{Tor}}}
\newcommand{\Div}{{\operatorname{Div}}}
\newcommand{\PDiv}{{\operatorname{PDiv}}}

\newcommand{\Char}{{\operatorname{char}}}
\newcommand{\Frac}{{\operatorname{Frac}}}
\newcommand{\Frob}{{\operatorname{Frob}}}
\newcommand{\Proj}{{\operatorname{Proj}\,}}
\newcommand{\RMod}{{}_R\operatorname{Mod}}
\newcommand{\SExt}{{\mathcal{E}\emph{xt}}}
\newcommand{\SHom}{{\emph{Hom}}}
\newcommand{\Span}{{\operatorname{span}}}
\newcommand{\Spec}{{\operatorname{Spec}\,}}
\newcommand{\Supp}{{\operatorname{Supp}}}
\newcommand{\Tens}[1]{#1\otimes --}
\newcommand{\fHom}[1]{\operatorname{Hom}(#1, --)}

\newcommand{\mSpec}{{\operatorname{mSpec}}}
\newcommand{\Grass}{{\operatorname{Grass}}}


% connectors
% \renewcommand{\mid}{\big|}
\newcommand{\sse}{\subseteq}
\newcommand{\subg}{\leqslant}
\newcommand{\supg}{\geqslant}
\newcommand{\ssne}{\subsetneq}
\newcommand{\isosubg}{\lesssim}
\newcommand{\nsse}{\not\subseteq}
\newcommand{\nsubg}{\trianglelefteq}
\newcommand{\nsupg}{\trianglerighteq}
\newcommand{\csubg}{\blacktriangleleft}
\renewcommand{\mid}{\big|}


% Arrows
\newcommand{\rw}{\rightarrow}
\newcommand{\Rw}{\Rightarrow}
\newcommand{\lw}{\leftarrow}
\newcommand{\Lw}{\Leftarrow}
\newcommand{\lrw}{\leftrightarrow}
\newcommand{\LRw}{\Leftrightarrow}
\newcommand{\acton}{\curvearrowright}
\newcommand{\actson}{\curvearrowright}
\newcommand\mapsfrom{\mathrel{\reflectbox{\ensuremath{\mapsto}}}}

% arrows for commutative diagram: create four new angled double-struck arrows
\newcommand\myrot[1]{\mathrel{\rotatebox[origin=c]{#1}{$\Rightarrow$}}}
\newcommand\NEarrow{\myrot{45}}
\newcommand\NWarrow{\myrot{135}}
\newcommand\SWarrow{\myrot{-135}}
\newcommand\SEarrow{\myrot{-45}}


%shorthands
\newcommand{\oln}[1]{{\overline{#1}}}
\renewcommand{\bf}[1]{\textbf{#1}}
\newcommand{\qn}{\quad\newline}


% limits
\newcommand{\Lim}{{\varprojlim}}
\newcommand{\coLim}{{\varinjlim}}
\newcommand{\Dlim}{\lim\limits_{\longrightarrow}}
\newcommand{\coDlim}{\lim\limits_{\longleftarrow}}
\newcommand{\cofHom}[1]{\operatorname{Hom}(--, #1)}


%for saturated ideals in the localization section (I don't like these, I think I'll delete)
\newcommand{\LIm}{{}^e }
\newcommand{\LPre}{{}^c }

%  ╔══════════════════════════════════════════════════════════╗
%  ║                    Custom Commands                       ║
%  ╚══════════════════════════════════════════════════════════╝

% swap varphi and phi
\let\temp\phi
\let\phi\varphi
\let\varphi\temp

\newcommand{\placeholder}{\_\_\_\_\_}

\newcommand{\set}[2]{\left\{#1 \ \middle|\ #2\right\}}

%mod should not have the crazy amount of space to its right!
\renewcommand{\mod}{\,\operatorname{mod}\,}

% dcups
\newcommand{\dcup}{\sqcup}
\newcommand{\Dcup}{\coprod}
% \newcommand{\Dcup}{\bigsqcup}

\newcommand{\nbot}{\mathrel{\rlap{$\bot$}\mkern2mu/}}

%a nice large circle
\newcommand{\tikzcircle}[2][red,fill=black]{\tikz[baseline=-0.5ex]\draw[#1,radius=#2] (0,0) circle ;}%

\makeatletter
\newcommand*\bigcdot{\mathpalette\bigcdot@{.5}}
\newcommand*\bigcdot@[2]{\mathbin{\vcenter{\hbox{\scalebox{#2}{$\m@th#1\bullet$}}}}}
\makeatother

%somehow not a default command
\def\rddots{\mathstrut^{.^{.^{.}}}}

%% new delimiter
\DeclarePairedDelimiter{\ceil}{\lceil}{\rceil}


%  ╔══════════════════════════════════════════════════════════╗
%  ║                     environments                         ║
%  ╚══════════════════════════════════════════════════════════╝


%remember the option: breakable

% create tcolor boxes
\newtcbtheorem[number within=section]{thm}{Theorem}%
{
colback=green!5,
colframe=green!35!black,
fonttitle=\bfseries,
label type=theorem
}{th}

\newtcbtheorem[number within=section, use counter from=thm]{lem}{Lemma}%
{
colback=blue!5,
colframe=black!35!black,
fonttitle=\bfseries,
label type=lemma
}{lm}
\newtcbtheorem[number within=section, use counter from=thm]{prop}{Proposition}%
{
colback=white!5,
colframe=white!35!black,
fonttitle=\bfseries,
label type=proposition
}{pr}
\newtcbtheorem[number within=section, use counter from=thm]{cor}{Corollary}%
{colback=blue!5,
colframe=blue!35!black,
fonttitle=\bfseries,
label type=corollary
}{co}
\newtcbtheorem[number within=section, use counter from=thm]{axiom}{Axiom}%
{colback=black!5,
colframe=yellow!35!black,
fonttitle=\bfseries,
label type=axiom
}{ax}
\newtcbtheorem[number within=section, use counter from=thm]{defn}{Definition}%
{colback=red!5,
colframe=red!35!black,
fonttitle=\bfseries,
label type=definition
}{df}


% for <s-cr> to create \item[$\LRw$] instead of \item
\newenvironment{equivEnumerate}
 {\begin{enumerate}
	}{
\end{enumerate}}

%insure the exercises are properly numbered
\counterwithin{Exercise}{section}
\counterwithin{Answer}{section}


%Custom enviroments
 \newenvironment{note}
 {\begin{tcolorbox}[enhanced, sharp corners, colback=white , borderline={0.2pt}{0pt}{black}]
   \begin{quote}\textbf{Note}:
   }{
   \end{quote}\end{tcolorbox}}

 \newenvironment{intuition}
 {\begin{quote}\textbf{Intuition}:
   }{
   \end{quote}}

\newtcbtheorem[number within=section, use counter from=thm]{titledBox}{}{%
  enhanced,
  sharp corners,
  colback=white,
  colframe=black,
  borderline={0.2pt}{0pt}{black},
  left=2mm,
  right=2mm,
  top=2mm,
  bottom=2mm,
  title style={opacity=1},
  attach title to upper,
  before upper={\textbf{\tcbtitletext}\quad},
  coltitle=black,
  fonttitle=\bfseries,
  separator sign={\quad}
}{box}


%better example and proof environments
\def\exampletext{Example} % If English
\colorlet{colexam}{red!55!black} % Global example color


\newtcbtheorem[number within=section, use counter from=thm]{example}{Example}{% \newtcolorbox[use counter=testexample]{testexamplebox}{%
% Example Frame Start
empty,% Empty previously set parameters
title={\exampletext \thetcbcounter #1},% use \thetcbcounter to access the testexample counter text
% Attaching a box requires an overlay
attach boxed title to top left,
   % Ensures proper line breaking in longer titles
   minipage boxed title,
% (boxed title style requires an overlay)
boxed title style={empty,size=minimal,toprule=0pt,top=4pt,left=3mm,overlay={}},
coltitle=colexam,fonttitle=\bfseries,
before=\par\medskip\noindent,parbox=false,boxsep=0pt,left=3mm,right=0mm,top=2pt,breakable,pad at break=0mm,
   before upper=\csname @totalleftmargin\endcsname0pt, % Use instead of parbox=true. This ensures parskip is inherited by box.
% Handles box when it exists on one page only
overlay unbroken={\draw[colexam,line width=.5pt] ([xshift=-0pt]title.north west) -- ([xshift=-0pt]frame.south west); },
% Handles multipage box: first page
overlay first={\draw[colexam,line width=.5pt] ([xshift=-0pt]title.north west) -- ([xshift=-0pt]frame.south west); },
% Handles multipage box: middle page
overlay middle={\draw[colexam,line width=.5pt] ([xshift=-0pt]frame.north west) -- ([xshift=-0pt]frame.south west); },
% Handles multipage box: last page
overlay last={\draw[colexam,line width=.5pt] ([xshift=-0pt]frame.north west) -- ([xshift=-0pt]frame.south west); }%
}{ex}



\newtcolorbox{ProofTcb}[1][]{%
    empty,
    title={\emph{Proof}: #1},
    attach boxed title to top left,
    minipage boxed title,
    boxed title style={empty,size=minimal,toprule=0pt,top=4pt,left=3mm,overlay={}},
    coltitle=black!55!black,
    fonttitle=\bfseries,
    before=\par\medskip\noindent,
    parbox=false,
    boxsep=0pt,
    left=3mm,
    right=0mm,
    top=2pt,
    breakable,
    pad at break=0mm,
    before upper=\csname @totalleftmargin\endcsname0pt,
    width=\dimexpr\textwidth-3mm\relax,
    overlay unbroken={\draw[black!55!black,line width=.5pt] ([xshift=-0pt]title.north west) -- ([xshift=-0pt]frame.south west); },
    overlay first={\draw[black!55!black,line width=.5pt] ([xshift=-0pt]title.north west) -- ([xshift=-0pt]frame.south west); },
    overlay middle={\draw[black!55!black,line width=.5pt] ([xshift=-0pt]frame.north west) -- ([xshift=-0pt]frame.south west); },
    overlay last={\draw[black!55!black,line width=.5pt] ([xshift=-0pt]frame.north west) -- ([xshift=-0pt]frame.south west); },
    #1
}

\newenvironment{Proof}[1][]
  {\begin{ProofTcb}[#1]}
  {\end{ProofTcb}}


%  ╔══════════════════════════════════════════════════════════╗
%  ║                    for Dynkin diagrams                   ║
%  ╚══════════════════════════════════════════════════════════╝

\def\row#1/#2!{#1_{\IfStrEq{#2}{}{n}{#2}} & \dynkin{#1}{#2}\\}
\newcommand{\tble}[1]{
   \renewcommand*\do[1]{\row##1!}
   \[
      \begin{array}{ll}\docsvlist{#1}\end{array}
   \]
}

%  ╔══════════════════════════════════════════════════════════╗
%  ║                         links                            ║
%  ╚══════════════════════════════════════════════════════════╝

\definecolor{LinkColour}{HTML}{A64A3B} % favourite
% \definecolor{LinkColour}{HTML}{8C3D32} % 2nd place
\hypersetup{
  colorlinks,
  citecolor=black,
  filecolor=magenta,
  linkcolor=LinkColour,
  urlcolor=blue,
}



%  ╔══════════════════════════════════════════════════════════╗
%  ║                          misc                            ║
%  ╚══════════════════════════════════════════════════════════╝

\stackMath %to be able to stack text in math mode
\pgfplotsset{compat=1.18}
\makeindex

%import options: all, bibliography, book-formatting, booksSetting, customEnvironments, dynkin, hw-formatting, language-setup, newcommands, packages, preamble, proofs, settings, tcolorbox, test.svg,
\begin{document}
\maketitle
\tableofcontents
\newpage

\chapter{Real Numbers}
In this chapter, I will cover the construction of the real numbers

\section{Construction Of The Real Numbers}

The natural numbers can be thought as essentially being defined on the human intuition of counting \footnote{possible
that other animals have the same inuition, like humming birds or horses, but right now, I'll stick to human's since as
far as we know, humans are the only creatures that can conseptualize numbers beyond 20}. There are bigger sets of
numbers, most of them generated by fillinwg in gaps.


But what does it mean to fill in gaps? When you were first introduced to arithmetics, you were introduced to the basic
operations: $+, -, \times, \div, \sqrt{}, \ln$.


However, Not all these operations will always return a value: in $\N$, there is no solution to the equation $3 - 5$, since $-2 \notin
\N$ \footnote{Or more precicely, no $a \in \N$ can be used as a solution to $3 - 5 = a$}. The integers extend the
naturals making them closed under subtraction, that is
\[
	\forall a \in \Z\,, \exists b \in \Z \text{ such that } a + b = 0
\]
Usually, $b$ is denoted $-a$, so $a + (-a) = 0$. Conventionally, this is shorthanded to $a -a = 0$.


Similarly, $\Q$ can extends $\Z$ by making them closed under division, that is
\[
	\forall a \in \Q\,, \exists b \in \Q \text{ such that } ab = 1
\]
Usually, $b$ is denoted $a^{-1}$ or $\frac{1}{a}$, so $aa^{-1} = 0$ or $\frac{a}{a} = 1$.

One might think this covers all possible ``gaps'' which existed before. However, we
\begin{prop}{$\sqrt{2}$ is irrational}{sqrt2Irrational}
	There does not exist a $a \in \Q$ such that $a = \sqrt{2}$
\end{prop}

\begin{Proof}
	For the sake of contradiction, let's assume that $\sqrt{2}$ is rational. Let $a = \frac{p}{q}$ ($p,q \in \Z$) where $gcd(p,q) = 1$.
	If $gcd(p,q) \ne 1$, simply divide $\frac{p}{q}$ by $gcd(p,q)$. Since we're assuming $\sqrt{2}$ is rational:
	\[
		\sqrt{2} = \frac{p}{q}
	\]
	thus
	\[
		2 = \frac{p^2}{q^2} \Rw 2q^{2} = p^2
	\]
	From here, we'll take 2 cases:
	\begin{enumerate}
		\item[] \bf{$p$ is odd}: If $p$ is even, then $p^2$ is odd. However, no matter what $q$ is, $2q^{2}$ is even. Thus
			$p$ cannot be odd
		\item[] \bf{$p$ is even}: since $gcd(p,q) = 1$, $q$ cannot be even, since then $gcd(p,q) = 2$. Thus, $q$ is odd.
			Since $p$ is even, it's divisble by $2$. Thus $p^2$ is divisible by $4$. However, $2q^2$ is only divisible
			by $2$, meaning that $p^2$ and $2q^2$ must be different numbers, and hence $p$ cannot be even.
	\end{enumerate}
	However, all numbers in $\Z$ are either even or odd, thus we reached a contradiction.
\end{Proof}

Thus, $\Q$ is not closed under $\sqrt{}$. Furthermore, numbers like $\phi$ (the golden ration) is also not rational.  We can fill in all these
gaps by saying that any polynomial equation must have a solution. Thus
\[
	x^2 + 1 = 0 \Rw x = \sqrt{2} \qquad x^2 - x - 1 = 0 \Rw x = \phi
\]

However, this too has it's limits. In 1873, Euler proved that $e$ cannot be the solution of a polynomial equation.
Later, it was shown that $\pi, \ln(2)$, and $e^\pi$ each cannot be the solution to a polynomial equation. So, a further
class of numbers was invented: If we can devise an alrogirthm to determine a number, the number is \emph{computable}.
Essentially, all numbers you can think about are computable ($e, \pi$, etc.). However, there do actually exist
\emph{uncomputable} numbers, numbers that we know exist, but we cannot divise an alrogirthm to define it! An example of
this would be the Chaitin number. So, we need to find a way of defining an even bigger cateogry of numbers, one with
hopefully no gaps.

The way we'll do this is by defining a larger set, based off of $\Q$, which will represent all the gaps:
\begin{defn}{Dedikend Cut}{dedikendCut}
	A \bf{cut} in $\Q$ is a pair of subsets $A, B$ of $\Q$ st
	\begin{enumerate}
		\item $A\cup B = \Q$, $A \ne \emptyset$, $B \ne \emptyset$, $A\cap B = \emptyset$
		\item if $a \in A$ and $b \in B$, then $a < b$
		\item $A$ contains no largest element
	\end{enumerate}
	Let $x = A|B$ denote the [dedikend] cut
\end{defn}

\begin{example}{here}{ex2}
	\begin{enumerate}
		\item $\left\{ r \in \Q | r < 0 \right\} | \left\{ r \in \Q : r \ge 0 \right\}$. Notice how this split kind of
			looks like $(-\infty, 0) |[0,\infty)$
		\item $\left\{ r \in \Q | r \le 0\ \text{or}\ r^2 < 2\right\} | \left\{ r \in \Q : r^2 \ge 2\right\}$. Notice how this split kind of
			looks like

			$(-\infty, \sqrt{2}) |[\sqrt{2},\infty)$
		\item Note that $\pm \infty$ is not a Dedekind cut, since that would imply $\emptyset|\Q$, or $\Q|\emptyset$,
			but $A$ and $B$ cannot be empty.
	\end{enumerate}
\end{example}

Note that in most constructions, $A$ is the more interesting set, and $B$ contains the ``rest'' of the rationals.

We will give a name to this set of cuts:
\begin{defn}{Real Numbers}{realNumbers}
	The set of all Dedekind cuts of $\Q$ is called the \bf{real numbers}, denoted $\R$. A \bf{real number} is a cut in $\Q$.
\end{defn}

We will justify calling this set a set of numbers soon. We'll first show that this indeed contains the rationals, and
that it no longer contains ``gaps'' (meaning it contains all previous numbers we talked about, and more).

Notice that every rational number can be constructed using a Dedekind cut:
\begin{titledBox}{}
	For any $q \in \Q$, $A|B = \left\{ r \in \Q | r < q \right\}| \left\{ r in \Q | r \ge q \right\}$ is a Dedekind cut
\end{titledBox}

Thus, we can identify every $q \in \Q$ with a $q^* \in \R$. In this way, we have that $\Q^* \sse \R$.

The next question is are there still gaps? Is there a place we can cut $\R$ to split it? This would mean we would need
$A|B$, where $A\cup B = \R$, $A\cap B = \emptyset$, and $A, B \ne \emptyset$. The way we'll solve this is by first
translating the order relation given to us in the definition of a dedikend cut to the real numbers:

\begin{defn}{defining an order relation on $\R$}{orderRelationOnR}
	If $x = A|B$ and $y = C|D$ are cuts such that $A \sse C$, then $x$ is \bf{less than or equal to} $y$ and we write $y
	\le y$. If $A \sse C$ and $A \ne C$, then $x$ is \bf{less than} $Y$ and we write $x < y$.
\end{defn}

It can easily be proven that this is indeed an order relation.
\begin{note}
	The definition is essentially only reliant on the first set in each dedikend cut. This further emphasizes that it's
	the main set we care about. The 2nd set can really be thought of as the ``remainder'' set
\end{note}

With the notion of an order, we can now also bound sets/numbers:

\begin{defn}{Upper And Lower Bound}{upperAndLowerBound}
	An element $M \in \R$ is an \bf{upper bound} for a set $S \sse \R$ (or $S$ is bounded from above by $M$) if each $s \in S$ satisfies
	\[
		s \le M
	\]
	Similarly for \bf{lower bound}.
\end{defn}

\begin{defn}{Least Upper Bound and Greatest Lower Bound}{lubAndglb}
	A \emph{least upper bound} (or supremum) is an upper bound less than any other upper bound for $S$. We denote the least upper
	bound for $S$ as $\sup(S)$.
\end{defn}

\begin{example}{example 3}{ex3}
	$-1$ is the least upper bound for the negative integers, while $4$ is an upper bound.
\end{example}

Not all sets have a lest upper bound. For example, in $\Q$, there is no least upper bound for $\sqrt{2}$ (we will show
this when we prove $\Q$ is dense in $\R$). A set that contains all it's least upper bounds is called a \emph{complete}
set. This is the key property that $\R$ has:

\begin{thm}{$\R$ Is Complete}{rIsComplete}
	The set $\R$ is \bf{complete} in the sense that it satisifes the least upper bound property: If $S$ is a nonemtpy
	subset of $\R$ and is bounded above in $\R$, then there exists a least upper bound for $S$
\end{thm}

\begin{Proof}
	Let $\CC \sse \R$ be a nonempty subset of $\R$ that is bounded from above (so $\forall X|Y \in \CC$, $X|Y < A|B$ for some
	$A|B \in \R$, or written in shorthand $\forall x \in \CC, x < y$ for some $y \in \R$). Define the following dedikend
	cut:
	\begin{align*}
		C &= \left\{ a \in \Q \mid A|B \in \CC \text{ such that } a \in A \right\}\\
		D &= \text{ rest of $\Q$}
	\end{align*}
	Let $z = C|D$. It is easy to check that $z$ is indeed a cut. We need to show that $z \le y$ for any upper bound $y$.
	Indeed, let $y = C'|D'$ be any upper-bound for $z$. Then $C \sse C'$. Thus, $z \le y$, showing that $z$ is indeed
	the least upper bound
\end{Proof}

With this property, we can return to trying to make a dedikend cut of the real numbers. Let's say $\mathcal{A}
| \mathcal{B}$ is a cut of the real numbers. By definition, $\mathcal{A}$ contains no largest elemetn. furthermore, each
$b \in \mathcal{B}$ is an upper-bound for $\mathcal{A}$. Therefore, there exists a $y = \sup(\mathcal{A})$. Therefore,
every dedikend cut in $\R$ occurs at a real number!

\subsection[thinking or reals as numbers]{Thinking of $\R$ as numbers}

I mentioned earlier that $\R$ can be thought of as numbers. But what does it mean to be a number? The term ``number'' is
actually an archatic term nowadays. Structures labeled as ``numbers'' are those which were historically labeled as
numbers, since they were the ones studied the most. However, different types of numbers can have incredibly different
properties. For example $\N$ and $\R$ are both labeled as numbers, but $\N$ and $\R$ are very different from one another!
In mathematics, we have expanded this notion greatly in the domain of \emph{Algebra}. For our purposes, we want $\R$ to
have all the operation's we have learnt in arithmetics:
\[
	+\qquad -\qquad\times\qquad\div
\]
and that they have all the properties we have grown accustom to:
\begin{enumerate}
	\item $a+b = b+a$
	\item $a+(b+c) = (a+b)+c$
	\item $\exists 0 \in \R$ such that for every $a \in \R$, $0 + a = a$
	\item $\forall a \in \R$, $\exists -a \in \R$ such that $a - a = 0$
	\item $a\cdot b = b\cdot a$
	\item $a\cdot (b\cdot c) = (a\cdot b)\cdot c$
	\item $\exists 1 \in \R$ such that for every $a \in \R$, $1 \cdot a = a$
	\item $\forall a \in \R$, $\exists a^{-1} \in \R$ such that $a \cdot a^{-1} = 1$
	\item $a(b+c) = ab + ac$
	\item if $x \le y$ then $x + z \le y + z$ for any $z \in \R$
	\item \ [$\R$ is complete]
\end{enumerate}

We have just shown the 10th and 11th property, so we'll go over some of the others. I say some, because it becomes quite tedious
to prove all of them. They are all duable, but we will quickly start accumulating many very similar cases.


To define $+$ such that $x + y = E|F \in \R$ works as we are familiar with, we define:
\begin{align*}
	E &= \left\{ r \in \Q | \text{for some $a \in A$ and some $c \in C$, we have $r = a + c$} \right\}\\
	F &= \text{ rest of $\Q$}
\end{align*}

This operation is naturally well-defined. By definition, we also have that $x +y = y + x$.

With this same idea, we can define the zero cut $0*$, and show that $0* + x = x$ for every $x \in \R$. We can also show
associativity in the same way

To define $\times$, it becomes a bit harder, since there are cases when $x > 0*$ or $x < 0*$. Howver, besides that, it's
the same principle.


Also how we say \emph{the} real numbers because this construction is unique, making the word ``number'' a good use,
since we usually think of numbers being the usual thing.

\subsection[What about further extending the real numbers?]{What about further extending $\R$?}

comment on the hyperreals. They contain the reals. They do not form a metric space. they are an ordered topology though.
Also no unique hyperreals

\section{Cauchy and Convergent Sequences}

We will take a second look at completeness. We have shown that number systems smaller than the real numbers leave
some gaps. We tried showing those by what seemed like hand-wavy methods. I gave examples like golden ration, then $\pi$,
then the chaitin number, to keep showing there are gaps, but that's an inconsistent way of finding gaps,
especially for more general sets which we might want to prove there is not gap. Is there one way to test, one way to
find gaps, that a system must satisfy in order to test to see whether
a set has gaps? Furthermore, can we think of the concept of completeness beyond the idea of numbers? Are there more
general spaces for which the concept of completeness applies? This is what we cover in this section.

The main idea that was used to capture this idea is that if we try to approach any potential ``hole'' with a sequence of
numbers, that number \emph{must} exist in our set. For example,
\[
	f: \N \rw \Q \qquad x_n \in (\sqrt{2} - \frac{1}{n}, \sqrt{2} + \frac{1}{n})\cap \Q
\]
where $x_n$ is some rational number inside the ever smaller open interval. Then this sequence approach $\sqrt{2}$,
however $\sqrt{2} \notin \Q$, and it doesn't technically \emph{converge} since for a sequence to converge, it must
converge \emph{to} a number. But in $\Q$, $\sqrt{2}$ is \emph{not} a number, so it can't converge in it! We thus define
what it means for a sequence to get arbitrarily close to ``itself'' as $n \rw \infty$. \footnote{I will quickly point
out that though I used a computable function to find the hole, and basically all examples one can think of will also
find be computable, you can have uncomputable Cauchy functions, meaning numbers like the Chitin number are also
considered}.

Notice further that this new criterion for finding holes relies on sequences, in particular sequences that
converge to one point. This might actually mean that completness is a notion that can be defined on spaces where there
is a single point a sequence can approach. There are actually a couple more conditions that are also required; when
taken together they define something called a uniform space. It's in the framework of uniform spaces that we can define
something called \emph{uniform properties}, which are properties that are \emph{global} to the space. Completeness is
a uniform property, since the entire space is uniform. If you heard of uniform continuity, this also is a uniform
property (we will cover uniform continuity in section ref:HERE).

However, there is little advantage to build-up the formal construction of uniform spaces, since we would only use very
briefly here to how completeness can be extended in full generality, but then not use it again (maybe forever in your
mathematical career). Instead, we will prove how a subset of uniform spaces which we use much more often can be called
complete. That subset is the class of metric spaces!\footnote{Historically, completeness of metrics was studied first.
Later, when the concept was abstracted, uniform spaces were defined in terms of pseudometrics, and were found to be the
natural generalization for studying completeness, but they are rarely used for current lack of utility beyond their use
in generalization}.


To start off with, let's formally define cauchy sequences in a metric space:
\begin{defn}{Cauchy Sequence}{cauchSeq}
     If \((X,d)\) is a metric space, we
    say that a sequence \(\left( x_{n} \right)\) is a Cauchy sequence if for
    every \(\epsilon > 0\), there exists an \(N\mathbb{\in N}\) such that
    \(d\left( x_{n},\ x_{m} \right) < \epsilon\) whenever \(n,m > N\).
\end{defn}

Intuitively, this means that after we get far enough in the sequence, any two points will be arbitrarily close to each
other. Cauchy sequences are more general than convergent sequences:
\begin{example}{ample 1}{ex1}
	Take  $X =
	\mathbb{Q}$, with the standard metric. Take the sequence:

	\[x_{1} = 1,x_{n + 1} = 1 + \frac{1}{1 + x_{n}}.\]

	Notice that \(d\left( x_{n},\ x_{m} \right) < \epsilon\), and therefor
	it's Cauchy, but \(d\left( x_{n},\ x \right) < \epsilon\) is not true
	because it's a infinite continued fraction, that will equal to the
	golden ratio \(\phi\). Since \(\phi\notin\mathbb{Q}\), it doesn't
	cenverge to any \(x\). Another example taking advantage of different metrics is:

	\begin{center}
	  $((0,1], d_{Eu})$, $x_{n} = \frac{1}{n}$\\
	  $X = \mathbb{R}, d(x,y) = |e^{x} - e^{y}|, (x_{n}) = -n$
	\end{center}
\end{example}
However, all convergent sequences are cauchy:

\begin{prop}{convergent sequence $\Rw$ Cauchy sequence}{convergeSeqImpliesCauchSeq}
    if \((x_{n})\) is a convergent sequence in
    \((X,d)\), then \(\left( x_{n} \right)\) is Cauchy. (Note that being
    Cauchy does not mean it converges.)
\end{prop}

\begin{proof}
    Suppose \(\left( x_{n} \right) \rightarrow x\) and fix
    an \(\epsilon > 0\). Choose an \(N\mathbb{\in N}\) such that
    \(d\left( x_{n},\ x \right) < \frac{\epsilon}{2}\) if \(n > N\). This
    same N shows that \(\left( x_{n} \right)\) is Cauchy, for if \(m,n > N\)
    then

    \[d\left( x_{n},\ x_{m} \right) \leq d\left( x_{n},\ x \right) + d\left( x,x_{m} \right) < \frac{\epsilon}{2} + \frac{\epsilon}{2} = \epsilon\]

    Which suffices to prove it.
\end{proof}


We can ask when the converse holds? That is, when is it guaranteed that when approaching a point, that point will exist?
This essentially is the idea of completeness we formulated earlier said in a different way, since it says any time we
approach what we might suspect is a gap, we can guarantee there is an element there! To formalize this intuition:
\begin{thm}{$\R$ Is Complete Using Cauchy Sequences}{rIsCompleteUsingCauchySequences}
	$\R$ is \emph{complete} with respect to Cauchy sequences in the sense that if $(a_n)$ is a sequecne of real numbers
	which obeys a Cauchy condition then it converges to a limit in $\R$
\end{thm}

\begin{Proof}
	P.18-19. Much shorter proof than in 257. Basically showing first that it's bounded ,than appropriately constructing
	a set that keeps bounded the sequence from above as $n \rw \infty$, then take the $\sup$ of that set. Then prove
	that the cauchy sequence converges using the $\sup$ value.
\end{Proof}

This will become the de-facto way one thinks about completeness, since it generalizes the completeness proof presented
in \ref{th:rIsComplete} for arbitrary metric spaces.


Notice though that it works for metric spaces, and that the proof relied on the metric at hand. You can actually then
make cauchy sequecnes in spaces that seem to be complete, but are not (ex. look at the $e^x$ metric on $[0,a]$ we
defined earlier). You can even do more than that: you can have different metrics and construct different number systems!
For exmaple, If we let $p$ be a prime number, then let $d(x,y) = |x - y|_p$ where $|a|_p = \frac{1}{p^k}$ where $p^k$ is
the largest power of $p$ that divides $a$, then this will define a different metric that is \emph{not} the euclidean
metric. This will in fact induce a different order on $\Q$ than the usual order, that is, it will no longer be that
\[
	1/2 \le 1/5
\]
but HERE. This links us back to our results from earlier in this chapter, were the construction of $\R$ was intimitely
linked to the order of $\Q$.

Since it's a metric space, it will define cauchy sequences, and so we can compelte the space too! The completion
of $\Q$ through this metric is called labeled $\Q_p$ and is called the $p$-adic numbers!


I remember hearing that the real numbers are the unique comlete ordered field, but the $p$-adics have the $p$-adic
order, so I don't know how to combine the two.

\section[Further properties of real numbers]{Further properties of $\R$}

In this section, we'll introduce a couple more proeprties the real numbers grant us.
\subsection*{Density}

In topology, Given a space $X$, a set $S$ is \emph{dense} if $\oln{S} = X$. We can topologically prove that $\oln{\Q}
= \R$. In real analysis, we would always be taking subsets $S \sse \R$ such that $\oln{S} = \R$. Density in $\R$ has
a special property:
\begin{prop}{Density In $\R$}{densityInR}
	A subset $S \sse \R$ is said to be dense in $\R$ if between any two real numbers, there exists an element of $S$.
	That is, if for any real number $a,b$ such that $a < b$, then we have $S \cap(a,b) \ne \emptyset$
\end{prop}

\begin{Proof}
	Exercice.
\end{Proof}

\begin{example}{density}{density}
	The most famous such examples is that $\Q$ is dense in $\R$. A problem with doing.
\end{example}



\subsection*{Archemidean principle}


\begin{defn}{Archimedean Property}{archimedeanProperty}
	Given an ordered space $X$, if for every $x \in X$, there is an integer $n$ that is greater than $x$, then $X$ has
	the \emph{archemedean property}
\end{defn}

\begin{example}{Archimedea}{archimedian}
	$\R$ has the archemedean property. For any $x \in \R$, there exists a $n = \ceil{x}+1$
\end{example}


\subsection*{$\epsilon$-principle}

Finally, we'll label a fact that is very common, and you might've implictely internalized it alraedy doing
continuity, convergent sequence, and limit proofs:

\begin{thm}{$\epsilon$ Principle}{epsilonPrinciple}
	If $a,b$ are real numbers and if for each $\epsilon > 0$, we have $a \le b + \epsilon$, then $a \le b$. If for each
	$\epsilon > 0$, $|x-y| \le \epsilon$, then $x = y$
\end{thm}

\begin{Proof}
	easy to prove.
\end{Proof}

\chapter{Topologies Through the Lens of Analysis}

In this chapter, we will further elaborate on the standard topology of the real numbers given their special properties.

First, standard $\R$ is metrizable, so a reminder of the definition of metric from 327:

\begin{defn}{Metric}{metric}
	Let $X$ be a set. A \emph{metric} on $X$ is the function
	\[
		d: X \times X \rw \R_{\ge 0}
	\]
	such that
	\begin{enumerate}
		\item $d(x,y) = d(y,x)$
		\item $d(x,y) + d(y,z) \ge d(x,z)$
		\item $d(x,y) = 0 \text{ if and only if } x = y$
	\end{enumerate}
\end{defn}

\begin{note}
	the fact that it is an \emph{if and only if} is in fact very important. Without that you can say
	\[
		d(x,y) = 0 \qquad \forall x,y \in X
	\]
	and that would be a metric, and you can go on to say it would define the indiscreet topology, but this was purposely
	left out because it would then not make all metric spaces Hausdorff (as we'll show soon)
\end{note}
\begin{defn}{Metric Topology}{metricTop}
	If $d$ is a metric on $X$, we call the \emph{metric topology} on $X$, to e the topology generated by the basis
	\[
		\CB = \left\{ B_d(x,\epsilon) \mid x\in X, \epsilon > 0 \right\}
	\]
\end{defn}
Thus, standard $\R$ would have
\[
	d(x,y) = |x-y|
\]
\begin{defn}{Convergence in Metric Space}{convInMetricSp}
	A function $f: \N \rw M$ is said to \emph{converge} if
	\[
		\forall \epsilon > 0,\ \exists \N \in \N \text{ such that }\ \forall n \ge N, d(x_n, p) < \epsilon
	\]
\end{defn}

\begin{cor}{Convergence of sequence and sub-sequence}{convergenceOfSeqAndSubSeq}
	A sequence converges if and only if every subsequence converges. Furthermore, they converge to the same point
\end{cor}

\begin{Proof}
	Simple proof
\end{Proof}

\begin{defn}{Sequential Continuity}{sequentialCont}
	A function $f: M \rw N$ is \emph{continuous} if it preserves sequential convergence: $f$ sends convergent sequences
	in $M$ to convergent sequences in $N$, limits being sent to limits. That is, for each sequence $(p_n)$ in $M$ which
	converges to a limit $p$ in $M$, the image sequence $(f(p_n))$ converges to $fp$ in $N$.
\end{defn}

\begin{cor}{Continuity $\LRw$ sequential continuity in metric spaces}{contIsSeqContInMetSp}
	Let $M$ be a metric space. Then a function is continuous if and only if it is sequentially continuous.
\end{cor}

\begin{Proof}
	here
\end{Proof}

\begin{cor}{Epsilon-Delta Continuity}{epsilonDeltaCont}
	A function is continuous if
	\[
		\forall \epsilon > 0,\ \exists \delta > 0 \text{ such that }\ d(x,t) < \delta \Rw d(f(x), f(t)) < \epsilon
	\]
\end{cor}

\begin{Proof}
	here
\end{Proof}

\begin{defn}{Homeomorphim}{homeo}
	A function is homeomorphic if it's bijective and bi-continuous
\end{defn}

\begin{defn}{Closed Set In Metric Spaces}{closedSetInMetricSps}
	A set $S$ is a \bf{closed set} if it contains all its limit points
\end{defn}

This definition contrained to real-analysis is intereting, since it mirrors the idea of closure in different domains of
mathematics. In Algebra, an algebraically closed field is one for which every binary operation produces a number, in
particular $\sqrt{\_}$. Closed here gives a similar connotation.

\begin{defn}{Open Set In Metric Space}{openSetInMetricSp}
	A set is an \bf{open set} if for each $p \in S$ there exstis an $r > 0$ st
	\[
		d(p,q) < r \qquad \Rw \qquad q \in S
	\]
\end{defn}

We will use the following notaiton (where $M$ is the metric space)
%\begin{align*}
	%\lim S &= \left\{ p \in M | p \text{ is a limit of $S$} \right\} = \oln{S}\\
	%M_rS &= \left\{ p \in M | d(p,q ) < r\right\}
%\end{align*}

Note also the \emph{neigborhood} will refer to \emph{open neigbhorhoods}.

We also define a product metric:
\begin{defn}{Metric On Products Space}{metricOnProdsSp}
	The following are metric we usually define on $M = X\times Y$:
	\begin{align*}
		d_E(m, m') &= \sqrt{d_X(x,x')^2 + d_Y(y,y')^2}\\
		d_{max}(m, m') &= \max \left\{ d_X(x,x'), d_Y(y,y') \right\}\\
		d_{sum}(m, m') &= d_X(x,x') = d_Y(y,y')\\
	\end{align*}
\end{defn}

\begin{prop}{order-relation between metrics}{ordRelBetwnMetrics}
	\[
		d_{max} \le d_E \le d_{sum} \le 2d_{max}
	\]
\end{prop}

\begin{Proof}
	here
\end{Proof}

Then here there is all your knowledge about compactness

Then here there is nests of compactness
If $A_1 \supset A_2 \cdots \supset A_n \supset A_{n+1} \supset ... $ Then $(A_n)$ is a nested sequence of sets. Its
intersection is:
\[
	\cap_{i=1}^\infty A_n = \left\{ p | \text{for each $n$ we have $p \in A_n$} \right\}
\]

\begin{thm}{intersection of nested compact sequence}{intersectionOfNestCompSeq}
	The intersection of a nested sequence of compact nonemppty sets is compact and non-empty
\end{thm}

\begin{Proof}
	here
\end{Proof}

The diameter of a nonempty set $S \sse $ is the supremum of the distances $d(x,y)$ between points of $S$

\begin{cor}{Nested Compact Sequence to one-point}{nestCompSeqTo1Point}
	If in addition to being nested, nompety, and ompact, the sets $A_n$ have diameter that tends to $0$ as $n \rw
	\infty$ thne $A = \cap A_n$ is a single point
\end{cor}

\begin{Proof}
	p.83
\end{Proof}

Next is the concept of uniform continuity, which preserves global properties of spaces (ex. completness, compactness,
etc.)

\begin{defn}{Uniform Continuity}{uniformContinuity}
	here
\end{defn}

\subsection{Clustering and Perfec Metric Space}

\begin{defn}{Limit Point}{limitPoint}
	here
\end{defn}

\begin{defn}{Cluster Point}{clusterPoint}
	Let $M$ be a metric space and let $p \in M$. Then $p$ is a cluster point if for every $r > 0$
	\[
		|B_r(p)| \ge \aleph_0
	\]
	IF $S \sse M$ is a set such that $p$ is a cluster point around $S$, then we say $S$ \emph{clusters} at $p$.
\end{defn}

\begin{defn}{Condensation Point}{condensationPoint}
	Let $M$ be a metric space and let $p \in M$. Then $p$ is a condenses point if for every $r > 0$
	\[
		|B_r(p)| \ge \aleph_1
	\]
	IF $S \sse M$ is a set such that $p$ is a condensation point around $S$, then we say $S$ \emph{condenses} at $p$.

\end{defn}

Now, something linking metric spaces with completeness:


\begin{defn}{Perfect Metric}{perfectMetric}
	A metric space $M$ is \bf{perfect} if $M' = M$, i.e., each $p \in M$ is a cluster point of $M$. Recall that $M$
	clusters at $p$ if each $M_rp$ is an infinite set.
\end{defn}

\begin{example}{prefect metric}{perfectMetric}
	\begin{enumerate}
		\item $[a,b]$ is prefect and
		\item $\Q$ is perfect (meaning perfect $\not\Rw$ complete).
		\item $N \sse \R$ is not perfect since non of its points are cluster points
		\item If we generalize the notion of perfect to a more topological framework, a set is perfect if it's closed
			and has no isolated points (that is, $x$ is an isolated point of $S$ is there exists a neigborhood $U$ such
			that $x \in U$ and $U\cap S = \left\{ x \right\}$). Then $\N$ with the trivial topology has only perfect
			subsets, even though there are points not infinitely close to each other. This ``infintely close'' fact
			comes from every $\epsilon$-ball needing to not have isolated points (so it in a sense results from the fact
			that the co-domain of a metric is $\R$)
	\end{enumerate}
\end{example}

\begin{thm}{Nonempty, Perfect, Complete Metric Space Is Uncountable}{perfectAndUncountable}
 Every Nonempty, Perfect, Complete Metric Space Is Uncountable
\end{thm}

\begin{Proof}
	Using the fact that there must be at least an infinite number of ponts, we can at least create a countably long list
	of points. With this information try to construct a case to use Cantor's Diagonalization argument.
\end{Proof}


\begin{cor}{Locally Uncountable}{locallyUncountable}
	Every nonempty perfect complete metric space is everywehre uncountable in the sense that each $r$-neigborhood is
	uncountable.
\end{cor}

\begin{Proof}
	Take any $\epsilon$-ball around any point $x$. take $\frac{\epsilon}{2}$. Taking it's closure will make it remain in
	the $\epsilon$-ball. By \ref{th:perfectAndUncountable}, it is uncountable, and thus our $\epsilon$-ball is
	uncountable
\end{Proof}


\subsection{Arithmetical operations and continuity}

Fun exercice: verify $+, -, \cdot, \div$, and taking $1/x$ are all continuous (with $\div$ when the denominator is not
zero). So is the max, min, and $|x|$ functions. So is scalar multiplication. This immediately implies polynomials are
alway continuous.


\section{Compactness}


In metric spaces, we can define the concept of sequential compactness:
\begin{defn}{Sequentially Compact}{sequentiallyComp}
	A subset $A \sse X$ of a metric space $(X,d)$ is \emph{sequentially compact} if for every sequence in $A$, there
	exists a convergent subsequence that converges in $A$
\end{defn}

\subsection{Compactness in compact metric spaces}

\begin{thm}{Generalized Heine-Borel Theorem}{generalizedHeine-Borelorem}
	A subset of a complete metric space is compact if and only if it is closed and totally bounded.
\end{thm}

\begin{Proof}
	p.104 Pugh
\end{Proof}

Sometimes, the following corrollary is known as generalized Heine Borel. In $\R$:
\begin{center}
	Compact $\Rw$ closed and bounded
\end{center}
and for the converse in metric spaces:
\begin{center}
	Compact $\LRw$ complete and totally bounded
\end{center}
which is the result of his next corrollary
\begin{cor}{Compactness Of Metric Space (Generlized Heine-Borel)}{compnessMetricSp}
	A metric space is compact if and only if it is complete and totally bounded
\end{cor}

\begin{Proof}
	p.104
\end{Proof}

\subsection*{summary}
[covering] compactness : every open cover has finite subcover


sequentially compcat: every sequence has a convergent subsequence


In $\R$, Heine Borel provides:
\[
	\text{compact subset} \LRw \text{ closed and bounded subset }
\]
In metric space:
\[
	\text{compact subset } \Rw \text{ closed and bounded subset } \qquad \text{compact subset} \not\Lw \text{ closed and bounded subset}
\]
Or, if we think about it in terms of the entire space:
\[
	\text{compact space } \Rw \text{ bounded space } \qquad \text{compact space} \not\Lw \text{bounded space}
\]
counter example being $\Q \sse \R$
\begin{example}{Example}{ex4}
	Take $[0, \frac{1}{\pi} - \frac{1}{n}] \cup [ \frac{1}{\pi} + \frac{1}{n}, 1]$ for $n \ge 4$. This covers all of $\Q
	\cap [0,1]$, but has no fintie subcover. Note that these sets are open in the subspace topology, but not in $\R$.
\end{example}
However:
\[
	\text{compact space} \LRw \text{ complete and totally bounded }
\]
and os in complete metric space:
\[
	\text{compact space} \LRw \text{ totally bounded }
\]
and using theorem ref:HERE, we get:
\[
	\text{compact subset} \LRw \text{ closed and totally bounded subset}
\]


Note that $C^0([a,b], \R)$ is complete (a proof in prop ref:HERE). However, the infinite-dimensionality of it means we
need another condition for compactness: Using Arzela Ascoli, we get:
\[
	\text{compact subset} \LRw \text{ closed, uniformally bounded, and equicontinuous}
\]

For ``closed and bounded $\LRw$ compact'' in a metric space you need:
\[
	\text{complete, $\sigma$-compact, locally compact}
\]
which we'll never use in these notes, but it's a fun trivial.

\chapter{Differentiation}\label{cha:differentiation}

here

\chapter{Integration}\label{cha:integration}
here

\section{Differentiation Under the Integral Sign}

A fundamental question in analysis is identifying the conditions under which we can commute the integral and derivative operators. This is formalized by the Leibniz Integral Rule.

\begin{thm}{Leibniz Integral Rule (Constant Limits)}{leibnizRule}
    Let $f(x, t)$ be a function defined for $x \in [a, b]$ and $t \in [c, d]$. Suppose that:
    \begin{enumerate}
        \item $f(x, t)$ is continuous on the rectangle $[a, b] \times [c, d]$.
        \item The partial derivative $\frac{\partial f}{\partial t}$ exists and is continuous on $[a, b] \times [c, d]$.
    \end{enumerate}
    Then, $F(t) = \int_{a}^{b} f(x, t) \, dx$ is differentiable on $(c, d)$ and:
    \[
        \frac{d}{dt} \left( \int_{a}^{b} f(x, t) \, dx \right) = \int_{a}^{b} \frac{\partial}{\partial t} f(x, t) \, dx
    \]
\end{thm}

\begin{Proof}
    Let $F(t) = \int_{a}^{b} f(x, t) \, dx$. By the definition of the derivative:
    \[
        F'(t) = \lim_{h \to 0} \frac{F(t+h) - F(t)}{h} = \lim_{h \to 0} \frac{\int_{a}^{b} f(x, t+h) \, dx - \int_{a}^{b} f(x, t) \, dx}{h}
    \]
    By the linearity of the integral, we can combine the terms:
    \[
        F'(t) = \lim_{h \to 0} \int_{a}^{b} \frac{f(x, t+h) - f(x, t)}{h} \, dx
    \]
    We now fix $x$ and consider $f(x, t)$ strictly as a function of $t$. By the Mean Value Theorem, for any fixed $x$, there exists a value $c_{x,h}$ strictly between $t$ and $t+h$ such that:
    \[
        \frac{f(x, t+h) - f(x, t)}{h} = \frac{\partial f}{\partial t}(x, c_{x,h})
    \]
    We substitute this back into our expression. We wish to show that the limit of the integral approaches $\int \frac{\partial f}{\partial t}$. To do so, we analyze the difference:
    \[
        \left| \int_{a}^{b} \frac{\partial f}{\partial t}(x, c_{x,h}) \, dx - \int_{a}^{b} \frac{\partial f}{\partial t}(x, t) \, dx \right|
    \]
    Using the inequality $|\int g| \le \int |g|$:
    \[
        \le \int_{a}^{b} \left| \frac{\partial f}{\partial t}(x, c_{x,h}) - \frac{\partial f}{\partial t}(x, t) \right| \, dx
    \]
    Here we rely on the compactness of the domain. Since $\frac{\partial f}{\partial t}$ is continuous on the
    closed, bounded rectangle $[a, b] \times [c, d]$, it is \emph{uniformly continuous}.

    This implies that for any $\epsilon > 0$, there exists a $\delta > 0$ such that if the distance between two points is less than $\delta$, the difference in function values is less than $\frac{\epsilon}{b-a}$.

    Since $c_{x,h}$ is between $t$ and $t+h$, we know $|c_{x,h} - t| < |h|$. If we choose $|h| < \delta$, then for all $x \in [a, b]$:
    \[
        \left| \frac{\partial f}{\partial t}(x, c_{x,h}) - \frac{\partial f}{\partial t}(x, t) \right| < \frac{\epsilon}{b-a}
    \]
    Substituting this bound back into the integral:
    \[
        \int_{a}^{b} \left| \frac{\partial f}{\partial t}(x, c_{x,h}) - \frac{\partial f}{\partial t}(x, t) \right| \, dx < \int_{a}^{b} \frac{\epsilon}{b-a} \, dx = (b-a) \cdot \frac{\epsilon}{b-a} = \epsilon
    \]
    Thus, as $h \to 0$, the difference goes to 0, proving the equality.
\end{Proof}

\begin{note}
    If the limits of integration also depend on $t$, i.e., $a(t)$ and $b(t)$, and are differentiable, the general Leibniz rule applies:
    \[
        \frac{d}{dt} \int_{a(t)}^{b(t)} f(x, t) \, dx = f(b(t), t) \cdot b'(t) - f(a(t), t) \cdot a'(t) + \int_{a(t)}^{b(t)} \frac{\partial f}{\partial t} \, dx
    \]
\end{note}

\newpage
\printindex

\newpage
\printbibliography
\end{document}
